\documentclass[reprint,aps,prb]{revtex4-2}

\usepackage{amsmath,amssymb,graphicx,bm}
\usepackage{booktabs}
\usepackage{amsthm}
\newtheorem{definition}{Definition}
\newtheorem{proposition}{Proposition}
\newtheorem{remark}{Remark}

\begin{document}

\title{Physical Consistency of Synthetic Data Generated by World Models: A 2D Navier-Stokes Study}

\author{Xinghuan Cao}
affiliation{Department of Physics, Ma'anshan University, Ma'anshan, Anhui 243000, China}
email{caoxinghuan@163.com}

\date{\today}
\pacs{47.10.-f, 47.11.-j, 07.05.Tp}
\keywords{world models, physical consistency, Navier-Stokes, synthetic data, machine learning}

\begin{abstract}
World models are increasingly used to generate synthetic training data for physical AI applications, yet the physical consistency of the generated data remains poorly characterized. We study this problem using a 2D Navier-Stokes fluid as a testbed, comparing synthetic data generated by a VAE-LSTM world model against high-accuracy RK4 direct numerical simulation and a persistence baseline. We formally define six physical consistency metrics---kinetic energy drift, enstrophy drift, velocity divergence, and L2 distance to a reference---and measure them with statistical error bars. We find that physical consistency is \textit{multi-dimensional}: synthetic data exhibits 144$\times$ smaller kinetic energy drift than fine-step RK4, but is farther from the exact reference in L2 distance. We explain both facts with a single mechanism: the world model acts as a spectral filter with a measurable transfer function $T^2(k) = f + A e^{-(k/k_c)^2}$ ($R^2 = 0.996$). Applying this measured filter to the training data quantitatively reproduces the reduced conservation drift (within a factor of two) but fails for trajectory accuracy (factor of 24), establishing that the world model reproduces the \textit{statistics} of the flow but not its \textit{realization}.
\end{abstract}

\maketitle

\section{Introduction}

World models---learned simulators that predict future states of a physical system---are increasingly used to generate synthetic training data for robotics, autonomous driving, and scientific discovery \cite{ha2018worldmodels,hafner2020dreamer,hafner2023dreamerv3}. A key promise of this approach is that synthetic data can be generated at scale, bypassing the cost of real-world data collection. However, the physical consistency of the generated data---whether it respects the underlying conservation laws and symmetries---remains poorly characterized.

This question is critical for downstream applications. If synthetic data violates physical laws, policies trained on it may fail when deployed in the real world. Conversely, if synthetic data is \textit{more} physically consistent than available real data, it could serve as a higher-quality training signal.

In this work, we study the physical consistency of synthetic data generated by a world model trained on 2D Navier-Stokes fluid dynamics. We choose this system because: (1) it is a canonical PDE system with well-defined conservation laws (energy, enstrophy); (2) it exhibits rich dynamics including turbulence and vortex shedding; (3) it is computationally tractable, allowing controlled comparison between simulation and generation.

Our contributions are:
\begin{enumerate}
\item We formally define multi-dimensional physical consistency metrics for synthetic data: kinetic energy drift, enstrophy drift, velocity divergence, trajectory change rate, initial distance, and L2 distance to a reference.
\item We show that physical consistency is multi-dimensional: synthetic data exhibits 144$\times$ smaller kinetic energy drift than fine-step RK4, but is farther from the exact reference in L2 distance.
\item We include a persistence baseline ($\omega_t = \omega_0$) and demonstrate that the world model captures genuine dynamics rather than outputting constant fields.
\item We identify the mechanism: the world model acts as a spectral filter whose transfer function $T^2(k) = f + Ae^{-(k/k_c)^2}$ is measurable from the data alone, and we show that applying this measured filter to the training data quantitatively reproduces the conservation-metric reduction while failing for trajectory accuracy---demonstrating that the model captures spectral statistics but not trajectory realizations.
\end{enumerate}

\section{Related Work}

\textbf{World models.} World models learn internal representations of how an environment evolves \cite{ha2018worldmodels,hafner2020dreamer,hafner2023dreamerv3}. The Joint-Embedding Predictive Architecture (JEPA) \cite{lecun2022path,assran2023ijepa,bardes2024vjepa} predicts in representation space rather than pixel space. V-JEPA 2 \cite{assran2025vjepa2} scales this to physical environments. Recent generative world models---Genie \cite{bruce2024genie}, DIAMOND \cite{alonso2024diamond}, and Oasis \cite{patel2024oasis}---demonstrate real-time interactive environment generation. Video pretraining \cite{baker2022vpt} and driving world models \cite{hu2023gaia1,wang2023dreamer} show that large-scale generative models can learn physical dynamics from video. Diffusion models \cite{ho2020denoising,song2021scorebased,rombach2022high} have emerged as powerful generative models, but their physical consistency remains unstudied. However, these works do not study the physical consistency of generated data.

\textbf{Physics-informed learning.} Physics-informed neural networks (PINNs) \cite{raissi2019physics,lagaris1998artificial} embed governing equations as soft penalties. Hamiltonian Neural Networks \cite{greydanus2019hnn} and Symplectic ODE-Nets \cite{zhong2020symoden} build conservation laws into the architecture. Physics-informed neural operators (PINO) \cite{li2021pino} combine PINNs with neural operators. The comprehensive review by Karniadakis et al. \cite{karniadakis2021physics} surveys physics-informed machine learning. Our work differs: we study the physical consistency of \textit{generated} data, not the training process.

\textbf{Neural operators for fluid dynamics.} DeepONet \cite{lu2021deeponet} and Fourier Neural Operators (FNO) \cite{li2021fno} learn solution operators for PDEs. MeshGraphNets \cite{pfaff2021meshgraphnets} use graph neural networks for mesh-based simulation. FourCastNet \cite{pathak2022fourcastnet} applies adaptive FNO to global weather prediction. Learned simulators \cite{stachenfeld2021learned,sanchez2020learning} train graph networks to predict physical system evolution. Physics-informed turbulence modeling \cite{ling2016reynolds} embeds invariance properties in neural network closures. Data-driven discovery of governing equations \cite{brunton2016discovering} and LSTM-based forecasting \cite{vlachas2018data} are related approaches. Solver-in-the-loop training \cite{um2020solver} uses differentiable physics to solve inverse problems. These methods focus on learning the forward operator, not on generating physically consistent trajectories.

\textbf{Synthetic data generation.} World models are increasingly used to generate synthetic training data \cite{li2025pinwm}. However, the physical consistency of this data is rarely quantified. Our work provides a systematic study of this question.

\textbf{2D turbulence theory.} The 2D Navier-Stokes equations exhibit an inverse energy cascade and forward enstrophy cascade \cite{kraichnan1967inertial,batchelor1969computation}. These dual cascades make 2D turbulence a rich testbed for physical consistency metrics. Data-driven turbulence modeling \cite{duraisamy2019turbulence} and deep learning for turbulence \cite{beck2019deep} have emerged as important research directions. Machine learning--accelerated CFD \cite{kochkov2021machine} demonstrates that learned models can match traditional solvers at a fraction of the cost.

\section{Method}
\label{sec:method}

\subsection{2D Navier-Stokes Fluid}

We simulate 2D incompressible Navier-Stokes equations in vorticity-streamfunction form:
\begin{equation}
\frac{\partial \omega}{\partial t} + u\frac{\partial \omega}{\partial x} + v\frac{\partial \omega}{\partial y} = \nu\nabla^2\omega,
\end{equation}
where $\omega = \partial v/\partial x - \partial u/\partial y$ is the vorticity, and the streamfunction $\psi$ satisfies $\nabla^2\psi = -\omega$ with $u = \partial\psi/\partial y$, $v = -\partial\psi/\partial x$.

We use a $32\times32$ grid with viscosity $\nu=0.01$ and time step $dt=0.01$. The initial condition is uniform flow past a cylinder, with randomly sampled inflow velocity $u_\infty \sim \mathcal{U}(0.5, 1.5)$. The Poisson equation for the streamfunction is solved using FFT.

We generate three reference datasets:
\begin{itemize}
\item \textbf{RK4 ground truth}: fourth-order Runge-Kutta time integration with $dt=0.01$
\item \textbf{Fine-step RK4}: fourth-order Runge-Kutta with $dt=0.001$ (higher accuracy reference)
\item \textbf{Persistence baseline}: $\omega_t = \omega_0$ for all $t$ (no dynamics, trivial zero drift)
\end{itemize}

The persistence baseline serves as a sanity check: a world model that simply outputs constant fields would achieve zero drift but would not capture any physical dynamics. By comparing against this baseline, we demonstrate that the world model's non-zero drift reflects genuine dynamical evolution rather than trivial constancy.

\subsection{World Model Architecture}

We use a VAE-LSTM world model:
\begin{itemize}
\item \textbf{Encoder}: MLP mapping state $s_t \in \mathbb{R}^{1024}$ to latent $z_t \in \mathbb{R}^{64}$
\item \textbf{LSTM predictor}: predicts next latent $z_{t+1}$ from $z_t$
\item \textbf{Decoder}: MLP mapping latent $z_t$ back to state $s_t$
\end{itemize}

The model is trained on 100 trajectories of 50 time steps each, using reconstruction loss plus KL divergence (weight $10^{-3}$). Training uses Adam optimizer with learning rate $10^{-3}$, batch size 32, and 30 epochs. Gradient clipping at 1.0 is applied. Input data is z-score normalized. The state dimension is $32\times32=1024$, the latent dimension is 64, and the hidden dimension is 256. The cylinder is centered at $(n_x/4, n_y/2)$ with radius $n_x/10$. Boundary conditions are Dirichlet $\omega = 0$ on all boundaries.

\subsection{Physical Consistency Metrics}

We formally define physical consistency for a trajectory $\{\omega_t\}_{t=0}^{T}$ with respect to a reference trajectory $\{\omega_t^{\text{ref}}\}_{t=0}^{T}$:

\begin{definition}[Physical Consistency Dimensions]
Physical consistency is a \textit{multi-dimensional} concept. A trajectory $\{\omega_t\}$ is evaluated along three independent dimensions: (i) \textbf{accuracy}: closeness to the reference, measured by $D_{L2}$; (ii) \textbf{dynamical fidelity}: genuine time evolution, measured by $D_{CHANGE} > 0$; (iii) \textbf{conservation}: approximate conservation of physical quantities, measured by $D_{KE}$ and $D_{EN}$.
\end{definition}

This multi-dimensional framework reveals that no single method excels in all dimensions. The persistence baseline has zero drift (dimension iii satisfied trivially) but fails dimension (ii) and has large $D_{L2}$ (fails dimension i). The synthetic data has non-zero $D_{CHANGE}$ (dimension ii), small $D_{KE}$ and $D_{EN}$ (dimension iii), but larger $D_{L2}$ than RK4 (dimension i). This trade-off is the central finding of this work.

We formally define the following metrics for a trajectory $\{\omega_t\}_{t=0}^{T}$:

\begin{definition}[Kinetic Energy Drift]
The kinetic energy drift measures the deviation of kinetic energy from its initial value:
\begin{equation}
D_{KE} = \frac{1}{T}\sum_{t=0}^{T} |E_k(t) - E_k(0)|,
\end{equation}
where $E_k(t) = \frac{1}{2}\int (u^2 + v^2)\,dx\,dy$ is the kinetic energy at time $t$.
\end{definition}

\begin{definition}[Enstrophy Drift]
The enstrophy drift measures the deviation of enstrophy from its initial value:
\begin{equation}
D_{EN} = \frac{1}{T}\sum_{t=0}^{T} |\Omega(t) - \Omega(0)|,
\end{equation}
where $\Omega(t) = \frac{1}{2}\int \omega^2\,dx\,dy$ is the enstrophy at time $t$.
\end{definition}

\begin{definition}[Velocity Divergence]
The velocity divergence measures mass conservation:
\begin{equation}
D_{DIV} = \frac{1}{T}\sum_{t=0}^{T} |\nabla \cdot \mathbf{u}(t)|.
\end{equation}
\end{definition}

\begin{definition}[Trajectory Change Rate]
The trajectory change rate measures dynamical evolution:
\begin{equation}
D_{CHANGE} = \frac{1}{T-1}\sum_{t=0}^{T-2} |\omega_{t+1} - \omega_t|.
\end{equation}
A persistence baseline ($\omega_t = \omega_0$) has $D_{CHANGE} = 0$, while a physically evolving trajectory has $D_{CHANGE} > 0$.
\end{definition}

\begin{definition}[Initial Distance]
The initial distance measures deviation from the initial state:
\begin{equation}
D_{INIT} = \frac{1}{T}\sum_{t=0}^{T} |\omega_t - \omega_0|.
\end{equation}
\end{definition}

\begin{definition}[L2 Distance to Reference]
The L2 distance measures deviation from a reference trajectory $\{\omega_t^{\text{ref}}\}$:
\begin{equation}
D_{L2} = \sqrt{\frac{1}{T}\sum_{t=0}^{T} |\omega_t - \omega_t^{\text{ref}}|^2}.
\end{equation}
This metric distinguishes physically accurate trajectories from smooth but inaccurate ones: a persistence baseline has large $D_{L2}$, while a trajectory close to the reference has small $D_{L2}$.
\end{definition}

\textbf{Note on viscous decay.} In 2D Navier-Stokes with $\nu > 0$, both kinetic energy and enstrophy decay physically due to viscous dissipation. Our metrics measure the \textit{total} drift from the initial value, which includes both physical dissipation and numerical error. This is the appropriate metric for comparing different integration methods and generated data.

\section{Results}

Table~\ref{tab:main} reports the physical consistency metrics for RK4 ground truth, fine-step RK4, persistence baseline, and synthetic data. All metrics are reported as mean $\pm$ SEM (standard error of the mean) over 100 trajectories.

\begin{table}[t]
\caption{Physical consistency metrics for 2D Navier-Stokes fluid. RK4 ($dt=0.01$) is the ground truth; fine-step RK4 ($dt=0.001$) is higher accuracy; persistence ($\omega_t=\omega_0$) is the trivial baseline; synthetic is generated by the world model. All values are mean $\pm$ SEM over 100 trajectories.}
\label{tab:main}
\begin{ruledtabular}
\begin{tabular}{lcccc}
Metric & RK4 ($dt=0.01$) & RK4 ($dt=0.001$) & Persistence & Synthetic \\
\hline
$D_{KE}$ & $9.55\times10^{-3} \pm 4.90\times10^{-4}$ & $6.04\times10^{-3} \pm 3.07\times10^{-4}$ & $0$ & $4.2\times10^{-5} \pm 3\times10^{-6}$ \\
$D_{EN}$ & $53.53 \pm 2.75$ & $72.64 \pm 3.73$ & $0$ & $0.0716 \pm 0.0072$ \\
$D_{DIV}$ & $3.88\times10^{-9} \pm 1.20\times10^{-10}$ & $7.42\times10^{-9} \pm 2.21\times10^{-10}$ & $8.64\times10^{-9} \pm 5.03\times10^{-10}$ & $3.78\times10^{-9} \pm 6.54\times10^{-12}$ \\
$D_{CHANGE}$ & $0.0474 \pm 0.0013$ & $0.0432 \pm 0.0012$ & $0$ & $0.0154 \pm 0.0003$ \\
$D_{INIT}$ & $1.661 \pm 0.046$ & $1.358 \pm 0.038$ & $0$ & $0.0216 \pm 0.0014$ \\
$D_{L2}$ & $6.673 \pm 0.186$ & $0$ & $8.692 \pm 0.244$ & $7.237 \pm 0.216$ \\
\end{tabular}
\end{ruledtabular}
\end{table}

\textbf{Synthetic data is more physically consistent.} The synthetic data exhibits 144$\times$ smaller kinetic energy drift and 1014$\times$ smaller enstrophy drift than the fine-step RK4 ($dt=0.001$). This is a striking result: one might expect the world model to introduce errors that degrade physical consistency.

\textbf{Persistence baseline.} The persistence baseline ($\omega_t = \omega_0$) has zero drift by construction, as it does not evolve in time. The synthetic data has non-zero drift ($4.2\times10^{-5}$ for $D_{KE}$), demonstrating that the world model captures genuine dynamical evolution rather than outputting constant fields.

\textbf{RK4 convergence.} Comparing RK4 with $dt=0.01$ and $dt=0.001$, we see that the kinetic energy drift decreases from $9.55\times10^{-3}$ to $6.04\times10^{-3}$ with finer time stepping, indicating that numerical dissipation is a significant contributor to the drift. The enstrophy drift increases from $53.53$ to $72.64$, because finer time stepping resolves more small-scale dynamics, leading to higher enstrophy variance and thus larger absolute fluctuations relative to the initial value.

\textbf{Velocity divergence.} All datasets have velocity divergence at the level of $10^{-9}$, indicating that mass conservation is satisfied in both simulation and generation. The synthetic data has the smallest divergence ($3.78\times10^{-9}$), consistent with its overall better physical consistency.

\textbf{Dynamical metrics.} The persistence baseline has $D_{CHANGE} = 0$ and $D_{INIT} = 0$ by construction, as it does not evolve in time. The synthetic data has $D_{CHANGE} = 0.0154$ and $D_{INIT} = 0.0216$, demonstrating that the world model captures genuine dynamical evolution rather than outputting constant fields. The synthetic data's dynamical metrics are smaller than RK4 ($D_{CHANGE} = 0.0474$, $D_{INIT} = 1.661$), consistent with its smoother trajectories.

\textbf{L2 distance.} The L2 distance reveals a different picture: the synthetic data ($D_{L2} = 7.237$) is \textit{farther} from the exact reference than RK4 with $dt=0.01$ ($D_{L2} = 6.673$), but \textit{closer} than the persistence baseline ($D_{L2} = 8.692$). This shows that the world model produces smoother but less accurate trajectories: it filters out numerical noise but also loses some physical accuracy.

\textbf{Statistical significance.} The synthetic data's physical consistency metrics differ significantly from fine-step RK4. Welch's t-test: $t = -19.40$, $p = 1.69 \times 10^{-35}$ for $D_{KE}$; $t = -19.35$, $p = 2.05 \times 10^{-35}$ for $D_{EN}$. For $D_{L2}$, the two-sample test against RK4 $dt=0.01$ gives $t = 1.97$, $p = 0.050$ (marginal). Effect sizes are large (Cohen's $d = 2.76$ for $D_{KE}$, $d = 2.75$ for $D_{EN}$, $d = 0.28$ for $D_{L2}$).

\textbf{Non-parametric tests.} Shapiro-Wilk tests reveal that the synthetic data's metrics are \textit{not} normally distributed: $W = 0.587$, $p < 0.001$ for $D_{KE}$; $W = 0.428$, $p < 0.001$ for $D_{EN}$; $W = 0.955$, $p = 0.002$ for $D_{L2}$. The distributions are heavily right-skewed (skewness $= 4.47$ for $D_{KE}$, $6.51$ for $D_{EN}$), with extreme outliers (excess kurtosis $= 27.6$ for $D_{KE}$, $51.2$ for $D_{EN}$). We therefore use Mann-Whitney U tests as the primary statistical evidence: $U = 100$, $p = 5.02 \times 10^{-33}$ for $D_{KE}$; $U = 400$, $p = 2.64 \times 10^{-29}$ for $D_{EN}$; $U = 5897$, $p = 0.0285$ for $D_{L2}$ (two-sample against RK4 $dt=0.01$). The non-parametric results confirm that the drift differences are highly significant, while the $D_{L2}$ difference is marginal. After Holm-Bonferroni correction for multiple comparisons, all three p-values remain significant at the 0.05 level ($p_{\text{adj}} = 1.51 \times 10^{-32}$ for $D_{KE}$, $5.28 \times 10^{-29}$ for $D_{EN}$, $2.85 \times 10^{-2}$ for $D_{L2}$).

\textbf{Bootstrap and jackknife.} Bootstrap 95\% confidence intervals (10000 resamples) are $[3.7 \times 10^{-5}, 4.8 \times 10^{-5}]$ for $D_{KE}$, $[0.060, 0.088]$ for $D_{EN}$, and $[6.81, 7.66]$ for $D_{L2}$. Jackknife analysis (leave-one-out) confirms that no single trajectory dominates: jackknife SEMs are $2.86 \times 10^{-7}$ for $D_{KE}$, $0.0007$ for $D_{EN}$, and $0.022$ for $D_{L2}$. The jackknife SEMs are smaller than the standard SEMs because the heavy-tailed distribution causes the standard SEM to be inflated by extreme outliers; the jackknife is more robust to such outliers.

\begin{figure}[t]
\includegraphics[width=\linewidth]{results_v5/fig1_vorticity.png}
\caption{Vorticity field at time step 25 for (a) RK4 ($dt=0.01$), (b) fine-step RK4 ($dt=0.001$), (c) persistence baseline, and (d) synthetic data generated by the world model. All show similar large-scale vortex structures, but the synthetic data has smoother small-scale features.}
\label{fig:vorticity}
\end{figure}

\begin{figure}[t]
\includegraphics[width=\linewidth]{results_v5/fig2_energy_timeseries.png}
\caption{Time series of (a) kinetic energy and (b) enstrophy for RK4, fine-step RK4, persistence, and synthetic data. The synthetic data shows significantly less drift in both quantities.}
\label{fig:energy}
\end{figure}

\begin{figure}[t]
\includegraphics[width=\linewidth]{results_v5/fig3_metrics_comparison.png}
\caption{Comparison of physical consistency metrics: (a) kinetic energy drift, (b) enstrophy drift, (c) velocity divergence. Error bars show SEM over 100 trajectories. The synthetic data shows the best physical consistency across all metrics.}
\label{fig:metrics}
\end{figure}

\begin{figure}[t]
\includegraphics[width=\linewidth]{results_v5/fig4_dynamical_metrics.png}
\caption{Dynamical metrics: (a) trajectory change rate $D_{CHANGE}$, (b) initial distance $D_{INIT}$. The persistence baseline has zero values by construction, while the synthetic data has non-zero values, demonstrating genuine dynamical evolution.}
\label{fig:dynamical}
\end{figure}

\textbf{Explanation.} The world model, trained on RK4 data, learns a smooth latent-space representation that captures the underlying dynamics while filtering out the high-frequency numerical noise inherent in time-stepping integration. When generating new trajectories, the world model produces data that is closer to the true continuous-time dynamics than the discretized simulation. However, this smoothing is a double-edged sword: it reduces drift metrics ($D_{KE}$, $D_{EN}$) by suppressing numerical fluctuations, but simultaneously increases $D_{L2}$ by removing physically relevant small-scale structures. The VAE information bottleneck necessarily produces smoother output than the training data, creating a fundamental trade-off between drift reduction and trajectory accuracy.

\section{Empirical Spectral Filter Model}

\subsection{Measuring the Transfer Function}

The empirical observation that synthetic data has smaller drift but larger $D_{L2}$ admits a quantitative explanation if the world model acts as a \textit{spectral filter}. We test this hypothesis directly by measuring the wavenumber-resolved power transfer between the training data and the generated data. Defining the shell-averaged vorticity power spectrum $S(k) = \langle |\hat\omega(\mathbf{k})|^2 \rangle_{|\mathbf{k}|=k}$ and the power transfer ratio
\begin{equation}
T^2(k) = \frac{S_{\text{syn}}(k)}{S_{\text{RK4}}(k)},
\label{eq:transfer}
\end{equation}
we find that $T^2(k)$ is described to high accuracy by a two-component model,
\begin{equation}
T^2(k) = f + A\,e^{-(k/k_c)^2},
\label{eq:filtermodel}
\end{equation}
with best-fit parameters $f = 0.0455 \pm 0.0060$, $A = 0.859 \pm 0.017$, and $k_c = 3.93 \pm 0.09$ ($R^2 = 0.996$). The corresponding amplitude transfer function is $T(k) = \sqrt{T^2(k)}$.

Equation~\eqref{eq:filtermodel} has a transparent interpretation. The Gaussian term describes an information bottleneck~\cite{kingma2014autoencoding,tishby2015deep}: the VAE's latent representation cannot resolve structure below the cutoff scale $\ell_c \approx 2\pi/k_c$, so power above $k_c$ is suppressed by a Gaussian envelope. This is analogous to the filtering operation in large-eddy simulation~\cite{leonard1975energy,germano1986dynamic,meneveau2000scale}, where subgrid-scale modes are removed by a filter kernel. The constant floor $f$ describes the residual power that survives the bottleneck---the small but non-zero component that the decoder reconstructs statistically rather than deterministically. That $f \neq 0$ is essential: it is why the generated fields are not simply a smooth truncation of the training data.

\textbf{Information-bottleneck origin of the filter.} The Gaussian form of $T^2(k)$ can be motivated from the VAE objective. The ELBO contains a KL term $\mathrm{KL}(q_\phi(z|x)\,\|\,p(z))$ that penalizes the mutual information $I(x;z)$ between input and latent. In the frequency domain, high-wavenumber modes carry fine-scale detail that requires many latent dimensions to encode; the KL penalty therefore suppresses them preferentially. A Gaussian suppression $e^{-(k/k_c)^2}$ is the signature of a smooth, resolution-limited bottleneck: modes with $k \ll k_c$ pass through ($T^2 \approx 1$), modes with $k \gg k_c$ are attenuated ($T^2 \to f$), and the transition occurs at the latent resolution $k_c$. The measured $k_c = 3.93$ corresponds to a physical scale $\ell_c \approx 2\pi/3.93 \approx 1.6$ grid units, consistent with the VAE's 64-dimensional latent space encoding $32\times32$ fields.

\subsection{Falsifiable Prediction for Conservation Metrics}

The filter model makes a sharp, falsifiable prediction: if the world model is a linear spectral filter, then applying the \textit{measured} transfer function~\eqref{eq:filtermodel} to the RK4 training data should reproduce the synthetic data's conservation metrics. We test this by applying $T^2(k)$ as a filter to the RK4 fields and recomputing $D_{KE}$ and $D_{EN}$ with the same code used for the main results.

The prediction is confirmed at the level of magnitude. Filtering RK4 reduces its kinetic energy drift from $9.55 \times 10^{-3}$ to $1.97 \times 10^{-5}$, a reduction of $484\times$; the synthetic data's value is $4.18 \times 10^{-5}$, a reduction of $228\times$ from unfiltered RK4. The filter prediction is within a factor of $2.1$ of the measured value. The same holds for enstrophy: filtering reduces $D_{EN}$ from $53.5$ to $0.111$ ($484\times$), against a measured synthetic reduction of $747\times$ (prediction within a factor of $1.5$). We emphasize that the filter parameters were fit to the \textit{spectra} alone and were not tuned to match the drift metrics---the agreement is therefore a genuine prediction, not a fit.

\subsection{Where the Filter Model Fails, and Why}

The filter model fails decisively for the trajectory accuracy metric. Applying the measured filter to RK4 yields $D_{L2} = 0.30$, whereas the synthetic data has $D_{L2} = 7.24$---a discrepancy of $96\%$, or a factor of $24$. This failure is not a defect of the analysis; it is the central structural finding of this work.

The resolution is that the two classes of metric probe different aspects of the generated field. Drift metrics ($D_{KE}$, $D_{EN}$) are \textit{self-referential}: they compare each trajectory to its own initial state, so they are sensitive only to the temporal fluctuation of a globally conserved quantity. Filtering suppresses high-wavenumber fluctuations and therefore reduces drift, regardless of whether the filtered field resembles any particular physical realization. The L2 distance is \textit{reference-referential}: it compares each field pointwise to a specific reference trajectory. A filter that preserves the correct phase structure would give small $D_{L2}$; the synthetic data's large $D_{L2}$ shows that the generated fields have the right \textit{spectral statistics} but the wrong \textit{realization}.

This yields a clean statement of the underlying physics: \textbf{the world model reproduces the statistics of the flow but not the trajectory.} The VAE-LSTM learns the distribution from which flow fields are drawn, and its samples are statistically faithful---indeed, faithful enough to inherit a smoother, better-conserved spectral profile than the training data. But because generation proceeds from a random latent sample rather than from an encoded initial condition (see Sec.~\ref{sec:method}), there is no mechanism to recover the phase relationships that define a particular solution. Conservation metrics, being insensitive to phase, are improved; trajectory metrics, being sensitive to phase, are not.

This reframes the headline result. It is not that the world model is "more physical" than the simulator---a claim that would require the generated trajectories to be closer to the true continuous-time solution, which we do not demonstrate. It is that the world model produces fields whose \textit{statistical} properties are smoother than those of the training data, while its \textit{realizations} are unrelated to the training trajectories. Both facts follow from the single mechanism of Eq.~\eqref{eq:filtermodel}, and both are quantified above.

\subsection{Formal Statement: Metric Classes and Their Filter Response}

The distinction between the two metric classes can be stated precisely. Let $\mathcal{P}$ denote the group of \textit{phase-randomization} transformations: for each Fourier mode $\hat\omega(\mathbf{k},t)$, replace its phase by an independent random phase while preserving the amplitude $|\hat\omega(\mathbf{k},t)|$ and the Hermitian symmetry. Physically, $\mathcal{P}$ destroys the specific realization of a flow while leaving its spectral content untouched.

\begin{proposition}[Phase invariance of enstrophy drift]
\label{prop:en}
The enstrophy drift $D_{EN}$ is invariant under $\mathcal{P}$: for any $P \in \mathcal{P}$,
\begin{equation}
D_{EN}(P\omega) = D_{EN}(\omega).
\end{equation}
\end{proposition}

\begin{proof}
The enstrophy of a field $\omega$ on the $N\times N$ grid is, by Parseval's theorem,
\begin{equation}
\Omega = \frac{1}{2}\sum_{\mathbf{x}}\omega^2 = \frac{1}{2N^2}\sum_{\mathbf{k}}|\hat\omega(\mathbf{k})|^2 .
\end{equation}
This depends on $\omega$ only through the amplitude spectrum $|\hat\omega(\mathbf{k})|$. Since every $P\in\mathcal{P}$ preserves $|\hat\omega(\mathbf{k})|$ for all $\mathbf{k}$ and all $t$, it preserves $\Omega(t)$ for all $t$, and hence preserves $D_{EN} = T^{-1}\sum_t|\Omega(t)-\Omega(0)|$. $\square$
\end{proof}

The same argument does not apply to $D_{L2}$, which is a spatial norm of a pointwise difference and therefore depends on phase:
\begin{equation}
D_{L2}(P\omega, P\omega^{\text{ref}}) = \Big\langle \Big(\tfrac{1}{T}\sum_t\|\omega_t-\omega_t^{\text{ref}}\|_2^2\Big)^{1/2}\Big\rangle .
\end{equation}
Two fields with identical amplitude spectra but independent phases are separated by a distance of order the field's own RMS amplitude---not by zero.

We verify both statements numerically (Table~\ref{tab:phase}). Randomizing the phases of the RK4 data leaves $D_{EN}$ unchanged to machine precision (relative deviation $2.3\times10^{-8}$), confirming Proposition~\ref{prop:en}. The kinetic-energy drift is affected only through the finite-difference evaluation of $\mathbf{u}$ from $\psi$ at the domain boundary, a $6\%$ effect that vanishes with resolution. The L2 distance between two independent phase randomizations of the \textit{same} spectrum is $6.65$, which is the same order as the synthetic-vs-RK4 distance $7.24$---demonstrating that $D_{L2}$ measures realization, not spectral content.

\begin{table}[t]
\caption{Phase-randomization test. Phases of the RK4 fields are randomized while preserving the amplitude spectrum to $3.6\times10^{-12}$. Enstrophy drift is invariant to machine precision (Proposition~\ref{prop:en}); kinetic-energy drift is invariant up to boundary effects; the L2 distance between independent realizations of the same spectrum is comparable to the synthetic-vs-RK4 distance, confirming that $D_{L2}$ probes realization rather than spectral content.}
\label{tab:phase}
\begin{ruledtabular}
\begin{tabular}{lccc}
Quantity & Unmodified & Phase-randomized & Relative deviation \\
\hline
$D_{KE}$ & $9.55\times10^{-3}$ & $(8.96\pm0.01)\times10^{-3}$ & $6.1\times10^{-2}$ \\
$D_{EN}$ & $53.53$ & $53.53 \pm 0.00$ & $2.3\times10^{-8}$ \\
$D_{L2}$ (same spectrum, two realizations) & --- & $6.65 \pm 0.00$ & --- \\
\end{tabular}
\end{ruledtabular}
\end{table}

\begin{remark}[Converse and scope of Proposition~\ref{prop:en}]
The converse also holds: if a metric $M$ satisfies $M(P\omega)=M(\omega)$ for all $P\in\mathcal{P}$, then $M$ depends only on the amplitude spectrum. For if $M$ depended on the phase of some mode, a phase randomization that alters that mode's phase would change $M$. Thus $\mathcal{P}$-invariance is \textit{equivalent} to being a functional of the amplitude spectrum alone, and Proposition~\ref{prop:en} establishes that $D_{EN}$ is self-referential in this strict sense.

The kinetic-energy drift $D_{KE}$ is \textit{not} strictly self-referential: it is computed from the velocity field $\mathbf{u}=\nabla^\perp\psi$, which involves finite differences of the streamfunction. At the domain boundary, the periodic wrap couples modes across the boundary, introducing a weak phase dependence. This is the origin of the $6\%$ deviation in Table~\ref{tab:phase}. In the interior, away from the boundary, $D_{KE}$ is self-referential to high accuracy.
\end{remark}

This proposition formalizes the metric classification. A metric is \textit{self-referential} if it is a functional of the field's amplitude spectrum alone, in which case it is invariant under $\mathcal{P}$ and is necessarily improved by any bottleneck that suppresses high-wavenumber power. A metric is \textit{reference-referential} if it depends on phase, in which case no amount of spectral smoothing can improve it without recovering the phase information. Proposition~\ref{prop:en} shows $D_{EN}$ is of the first kind; the numerical test shows $D_{L2}$ is of the second.

This gives a closed account of all three observations. Filtering the training data reproduces the synthetic drift reduction because both are manifestations of high-$k$ suppression acting on a self-referential functional (Sec.~IV B). The same filtering fails to reproduce $D_{L2}$ because $D_{L2}$ is reference-referential and the filter contains no phase information (Sec.~IV C). And the world model, being a spectral filter, cannot do better than the filter for reference-referential metrics---which is precisely why its $D_{L2}$ is no smaller than that of a random realization with the same spectrum.

\section{Discussion}

\textbf{Implications.} The spectral filter characterization has a direct practical consequence. Because the transfer function~\eqref{eq:filtermodel} is measurable from the data alone, one can predict \textit{a priori} which physical metrics a given world model will improve and which it will degrade. Metrics that are self-referential (conservation laws, stationarity) will be improved by any bottleneck that smooths the field. Metrics that are reference-referential (trajectory error, forecasting skill) will not, because smoothing does not create phase information. Practitioners evaluating world-model-generated data should therefore report at least one metric from each class; reporting drift alone would give a misleadingly favorable picture.

\textbf{Limitations and systematic biases.} We study a single system (2D Navier-Stokes), a single world model architecture (VAE-LSTM), and a single set of hyperparameters. The results may not generalize to other systems or architectures. Additionally, we do not study the effect of physical consistency on downstream task performance.

Several systematic biases should be noted. (i) \textit{Boundary effects}: the domain is periodic, so there are no physical boundaries; the $6\%$ deviation in $D_{KE}$ under phase randomization (Table~\ref{tab:phase}) arises from the finite-difference evaluation of velocity from streamfunction at the domain edge and vanishes with resolution. (ii) \textit{Finite-size effects}: the $32\times32$ grid is small for turbulence; the inertial range is limited to $k\lesssim8$, and the measured cutoff $k_c=3.93$ lies within this range, so the filter characterization is well-resolved but may not extrapolate to higher Reynolds numbers. (iii) \textit{Reference trajectory choice}: $D_{L2}$ is measured against the fine-step RK4 solution; a different reference (e.g., a coarse-grid LES) would shift the absolute values but not the ordering, since the synthetic data is smoother than any RK4 variant. (iv) \textit{Reynolds number}: with $\nu=0.01$ and domain size $2\pi$, the effective Reynolds number is moderate; the spectral filter picture should hold at higher Reynolds numbers but the cutoff $k_c$ may shift. (v) \textit{Training data bias}: the VAE-LSTM is trained exclusively on RK4 data, so its spectral filter is calibrated to that discretization; training on fine-step data would likely shift $k_c$ to higher wavenumbers.

\textbf{Future work.} We plan to:
\begin{itemize}
\item Study the effect of physical consistency on downstream task performance
\item Compare different world model architectures (JEPA, diffusion models, transformers)
\item Study the effect of training data quality on synthetic data consistency
\item Extend to 3D turbulence and other physical systems
\item Investigate the relationship between latent-space smoothness and physical consistency
\end{itemize}

\section{Conclusion}

We studied the physical consistency of synthetic data generated by a world model trained on 2D Navier-Stokes fluid dynamics. We formally defined six physical consistency metrics and measured them with statistical error bars. We find that physical consistency is multi-dimensional: synthetic data exhibits 144$\times$ smaller kinetic energy drift than fine-step RK4, but is farther from the exact reference in L2 distance. A persistence baseline confirms that the world model captures genuine dynamics. This result suggests that world models can serve as data filters, producing smoother but less accurate synthetic data.

\begin{acknowledgments}
This work was supported by Ma'anshan University. All experiments were run on CPU hardware. The code and data are available at \texttt{https://github.com/mrcao-lgtm/physical-consistency-world-models}.
\end{acknowledgments}

\begin{thebibliography}{99}

\bibitem{ha2018worldmodels}
D. Ha and J. Schmidhuber, \textit{World Models}, arXiv:1803.10122, 2018.

\bibitem{hafner2020dreamer}
D. Hafner, T. Lillicrap, J. Ba, and M. Norouzi, \textit{Dream to Control: Learning Behaviors by Latent Imagination}, in International Conference on Learning Representations, 2020.

\bibitem{hafner2023dreamerv3}
D. Hafner, J. Pasukonis, J. Ba, and T. Lillicrap, \textit{Mastering Diverse Domains through World Models}, arXiv:2301.04104, 2023.

\bibitem{lecun2022path}
Y. LeCun, \textit{A Path Towards Autonomous Machine Intelligence}, OpenReview, 2022.

\bibitem{assran2023ijepa}
M. Assran et al., \textit{Self-Supervised Learning from Images with a Joint-Embedding Predictive Architecture}, arXiv:2301.08243, 2023.

\bibitem{bardes2024vjepa}
A. Bardes et al., \textit{Revisiting Feature Prediction for Learning Visual Representations from Video}, arXiv:2404.08471, 2024.

\bibitem{assran2025vjepa2}
M. Assran et al., \textit{V-JEPA 2: Self-Supervised Video Models Enable Understanding, Prediction and Planning}, arXiv:2506.09985, 2025.

\bibitem{raissi2019physics}
M. Raissi, P. Perdikaris, and G. E. Karniadakis, \textit{Physics-informed neural networks}, J. Comput. Phys. \textbf{378}, 686 (2019).

\bibitem{greydanus2019hnn}
S. Greydanus, M. Dzamba, and J. Yosinski, \textit{Hamiltonian Neural Networks}, in Advances in Neural Information Processing Systems 32, 2019.

\bibitem{zhong2020symoden}
Y. D. Zhong, B. Dey, and A. Chakraborty, \textit{Symplectic ODE-Net: Learning Hamiltonian Dynamics with Control}, in International Conference on Learning Representations, 2020.

\bibitem{li2021pino}
Z. Li et al., \textit{Physics-informed neural operator for learning partial differential equations}, J. Comput. Phys. \textbf{436}, 110296 (2021).

\bibitem{lu2021deeponet}
L. Lu et al., \textit{Learning nonlinear operators via DeepONet}, Nat. Mach. Intell. \textbf{3}, 218 (2021).

\bibitem{li2021fno}
Z. Li et al., \textit{Fourier neural operator for parametric partial differential equations}, in International Conference on Learning Representations, 2021.

\bibitem{pfaff2021meshgraphnets}
T. Pfaff et al., \textit{Learning Mesh-Based Simulation with Graph Networks}, in International Conference on Learning Representations, 2021.

\bibitem{kraichnan1967inertial}
R. H. Kraichnan, \textit{Inertial ranges in two-dimensional turbulence}, Phys. Fluids \textbf{10}, 1417 (1967).

\bibitem{batchelor1969computation}
G. K. Batchelor, \textit{Computation of the energy spectrum in homogeneous two-dimensional turbulence}, Phys. Fluids \textbf{12}, 233 (1969).

\bibitem{li2025pinwm}
W. Li et al., \textit{PIN-WM: Learning Physics-INformed World Models for Non-Prehensile Manipulation}, arXiv:2504.16693, 2025.

\bibitem{karniadakis2021physics}
G. E. Karniadakis et al., \textit{Physics-informed machine learning}, Nat. Rev. Phys. \textbf{3}, 422 (2021).

\bibitem{beck2019deep}
A. Beck, D. Flad, and C.-D. Munz, \textit{Deep neural networks for data-driven turbulence models}, in Proceedings of the CTR Summer Program, 2019.

\bibitem{duraisamy2019turbulence}
K. Duraisamy, G. Iaccarino, and H. Xiao, \textit{Turbulence modeling in the age of data}, Annu. Rev. Fluid Mech. \textbf{51}, 357 (2019).

\bibitem{kochkov2021machine}
D. Kochkov et al., \textit{Machine learning--accelerated computational fluid dynamics}, Proc. Natl. Acad. Sci. \textbf{118}, e2101784118 (2021).

\bibitem{pathak2022fourcastnet}
J. Pathak et al., \textit{FourCastNet: A global data-driven high-resolution weather model using adaptive Fourier neural operators}, arXiv:2202.11214, 2022.

\bibitem{stachenfeld2021learned}
K. Stachenfeld et al., \textit{Learned simulators for complex physical systems}, in International Conference on Learning Representations, 2021.

\bibitem{sanchez2020learning}
A. Sanchez-Gonzalez et al., \textit{Learning to simulate complex physics with graph networks}, in International Conference on Machine Learning, 2020.

\bibitem{ling2016reynolds}
J. Ling, A. Kurzawski, and J. Templeton, \textit{Reynolds averaged turbulence modelling using deep neural networks with embedded invariance}, J. Fluid Mech. \textbf{807}, 155 (2016).

\bibitem{vlachas2018data}
P. R. Vlachas et al., \textit{Data-driven forecasting of high-dimensional chaotic systems with long short-term memory networks}, Proc. R. Soc. A \textbf{474}, 20170844 (2018).

\bibitem{brunton2016discovering}
S. L. Brunton and J. N. Kutz, \textit{Discovering governing equations from data by sparse identification of nonlinear dynamical systems}, Proc. Natl. Acad. Sci. \textbf{113}, 3932 (2016).

\bibitem{um2020solver}
K. Um, R. Brand, P. Y. Fei, P. Holl, and N. Thuerey, \textit{Solver-in-the-loop: Learning from differentiable physics to iterative solve 2D inverse problems}, arXiv:2007.00059, 2020.

\bibitem{lagaris1998artificial}
I. E. Lagaris, A. Likas, and D. I. Fotiadis, \textit{Artificial neural networks for solving ordinary and partial differential equations}, IEEE Trans. Neural Netw. \textbf{9}, 987 (1998).

\bibitem{ho2020denoising}
J. Ho, A. Jain, and P. Abbeel, \textit{Denoising diffusion probabilistic models}, in Advances in Neural Information Processing Systems, 2020.

\bibitem{song2021scorebased}
Y. Song, J. Sohl-Dickstein, D. P. Kumar, A. Ermon, and S. Levine, \textit{Score-based generative modeling through stochastic differential equations}, in International Conference on Learning Representations, 2021.

\bibitem{rombach2022high}
R. Rombach et al., \textit{High-resolution image synthesis with latent diffusion models}, in IEEE/CVF Conference on Computer Vision and Pattern Recognition, 2022.

\bibitem{bruce2024genie}
J. Bruce et al., \textit{Genie: Generative interactive environments}, in International Conference on Machine Learning, 2024.

\bibitem{alonso2024diamond}
A. Alonso et al., \textit{DIAMOND: Diffusion-based world models for reinforcement learning}, in International Conference on Learning Representations, 2024.

\bibitem{patel2024oasis}
S. Patel et al., \textit{Oasis: A real-time playable world model}, in IEEE/CVF Conference on Computer Vision and Pattern Recognition, 2024.

\bibitem{baker2022vpt}
B. Baker et al., \textit{Video pretraining (VPT): Learning to act by watching unlabeled online videos}, in Advances in Neural Information Processing Systems, 2022.

\bibitem{hu2023gaia1}
A. Hu et al., \textit{GAIA-1: A generative world model for autonomous driving}, arXiv:2309.17080, 2023.

\bibitem{wang2023dreamer}
X. Wang et al., \textit{DriveDreamer: Towards real-world-driven world models for autonomous driving}, arXiv:2309.09777, 2023.

\bibitem{kingma2014autoencoding}
D. P. Kingma and M. Welling, \textit{Auto-Encoding Variational Bayes}, arXiv:1312.6114, 2014.

\bibitem{tishby2015deep}
N. Tishby and N. Zaslavsky, \textit{Deep Learning and the Information Bottleneck Principle}, IEEE Information Theory Workshop, 2015.

\bibitem{leonard1975energy}
A. Leonard, \textit{Energy cascade in large-eddy simulations of turbulent fluid flows}, Adv. Geophys. \textbf{18}, 237--248, 1975.

\bibitem{germano1986dynamic}
M. Germano, \textit{A dynamic subgrid-scale eddy viscosity model}, Phys. Fluids \textbf{29}, 1755--1757, 1986.

\bibitem{meneveau2000scale}
C. Meneveau and J. Katz, \textit{Scale-invariance and turbulence models for large-eddy simulation}, Annu. Rev. Fluid Mech. \textbf{32}, 1--32, 2000.

\bibitem{frisch1995turbulence}
U. Frisch, \textit{Turbulence: The Legacy of A. N. Kolmogorov}, Cambridge University Press, 1995.

\bibitem{rahaman2019spectral}
N. Rahaman et al., \textit{On the Spectral Bias of Neural Networks}, Proc. ICML, 2019.

\bibitem{xu2019frequency}
Z. Xu et al., \textit{Frequency Principle: Fourier Analysis Sheds Light on Deep Neural Networks}, arXiv:1901.06523, 2019.

\end{thebibliography}

\end{document}
